% Copyright (c) 2026 Geovana Neves. Licensed under CC BY 4.0 (https://creativecommons.org/licenses/by/4.0/). See LICENSE-AND-AUTHORSHIP.md.
%
% overview/main.tex
%
% Guide 1 of the cfdbender guides: what a solver-agnostic workspace is, the
% spectrum of fidelities the library solvers cover, the workflow in one
% picture, a first campaign on the mock solver, and a map of guides 2 to 4.
% Built by ../build_guides.py with this folder and ../shared/ as include
% directories.
%
% No em dashes and no en dashes. No underscores in frame or section titles
% outside \code: beamer writes titles to .nav and .toc. A brace group right
% after a frame title is read as the subtitle and vanishes: a frame body that
% starts with a group starts with \relax.

\documentclass[aspectratio=169,11pt]{beamer}

\input{preamble-slides.tex}
\input{hooks-slides.tex}
\cfdbinput{boxes}
\cfdbinput{hooks-contract}

\usepackage[nohyphen]{underscore}

\input{info.tex}
\guidepdftitle{1}{Overview}

\newcolumntype{P}[1]{>{\RaggedRight\arraybackslash}p{#1}}
\newcolumntype{R}{>{\RaggedRight\arraybackslash}X}

\begin{document}

\guidetitleframe{1}{Overview}{Overview}
  {One workspace, many fidelities}
  {Sources: the package pages ``Workspace layout'', ``Command line'' and
   ``Matrix columns'' [1, 2, 3].}

\begin{frame}{Who reads what}
\begin{columns}[T]
\begin{column}{0.48\textwidth}
\begin{block}{This guide}
\begin{itemize}\small
  \item Part 1: what a solver-agnostic workspace is.
  \item Part 2: the solvers, as a spectrum of fidelities.
  \item Part 3: the workflow, the same for every solver.
  \item Part 4: a first campaign, with no solver installed.
  \item Part 5: the other guides and the order to read them in.
\end{itemize}
\end{block}
\end{column}
\begin{column}{0.48\textwidth}
\begin{block}{The other guides}
\begin{itemize}\small
  \item 2: workspaces. The tree, the inputs, plan, run, collect, post, a
        workstation and a cluster.
  \item 3: solvers. What each one reads and writes, what differs, how to
        choose.
  \item 4: the cheatsheet. Every command and option, and one page per stage
        of the workflow.
\end{itemize}
\end{block}
\end{column}
\end{columns}
\end{frame}

\begin{frame}{Contents}
\tableofcontents
\end{frame}

% Copyright (c) 2026 Geovana Neves. Licensed under CC BY 4.0 (https://creativecommons.org/licenses/by/4.0/). See LICENSE-AND-AUTHORSHIP.md.
%
% overview/text/01-idea-and-spectrum.tex
%
% Parts 1 and 2 of guide 1: the solver-agnostic workspace and the spectrum of
% fidelities.

\section{One workspace, many solvers}

\begin{frame}{The problem a workspace solves}
\begin{columns}[T]
\begin{column}{0.48\textwidth}
\begin{block}{Without it}
\begin{itemize}\small
  \item Each solver has its own input deck, its own folder habits and its
        own way of saying that a run ended.
  \item The study lives in scripts that only their writer can read.
  \item Changing solver means starting the bookkeeping again.
\end{itemize}
\end{block}
\end{column}
\begin{column}{0.48\textwidth}
\begin{block}{With it}
\begin{itemize}\small
  \item A study is \textbf{one folder}, the workspace, with the same tree
        for every solver.
  \item A study is \textbf{one table}, the run matrix: one row per
        simulation.
  \item The same commands plan it, run it, follow it and collect it.
\end{itemize}
\end{block}
\end{column}
\end{columns}
\begin{takeaway}You learn the workflow once. The solver changes what is
inside three input files, not how the campaign is run.\end{takeaway}
\end{frame}

\begin{frame}[fragile]{The solver is one line}
\begin{columns}[T]
\begin{column}{0.44\textwidth}
\begin{lstlisting}[style=cfdbtoml]
# init.toml, in the workspace root
schema_version = 1
solver = "su2"
\end{lstlisting}
\begin{itemize}\small
  \item \code{cfd-workspace init} writes it, once.
  \item Every command reads it first.
  \item A workspace keeps its solver for life.
\end{itemize}
\end{column}
\begin{column}{0.52\textwidth}
\relax{\cfdbtablesize
\begin{tabularx}{\linewidth}{@{}P{0.28\linewidth}R@{}}
\toprule
\textbf{Solver name} & \textbf{What runs} \\
\midrule
\code{su2} & SU2, by its executable or its Python wrapper \\
\code{openfoam} & OpenFOAM (ESI), \code{simpleFoam} \\
\code{pyfr} & PyFR, flux reconstruction \\
\code{xfoil} & XFOIL, airfoil polars \\
\code{qprop} & QPROP, an isolated propeller \\
\code{mock} & nothing: a stand-in to learn with \\
\code{flightstream} & FlightStream, handed to pyflightstream \\
\bottomrule
\end{tabularx}}
\end{column}
\end{columns}
\end{frame}

\begin{frame}{What stays the same, what changes}
\relax{\cfdbbodysize
\begin{tabularx}{\linewidth}{@{}P{0.16\linewidth}P{0.4\linewidth}R@{}}
\toprule
 & \textbf{The same for every solver} & \textbf{The solver's own} \\
\midrule
The tree & \code{inputs/}, \code{sims/}, \code{post/}, \code{archive/},
  \code{runs.json} & the files inside one point folder \\
The matrix & its columns and what they mean & a few flight-condition keys \\
The commands & every command and option of \code{cfd-matrix} and
  \code{cfd-workspace} & nothing \\
The inputs & one setup, one reference and one post file per row & the keys
  of the setup file, the mesh format \\
The result & status, verdict, one file set per polar & the coefficient names,
  what ``converged'' rests on \\
\bottomrule
\end{tabularx}}
\vspace{2mm}
\begin{takeaway}Guide 2 teaches the left column. Guide 3 teaches the right
one, solver by solver.\end{takeaway}
\end{frame}

\section{A spectrum of fidelities}

\begin{frame}{More than one wind tunnel}
\begin{itemize}
  \item A laboratory runs more than one wind tunnel: a small one answers many
        cheap questions, a large one answers a few expensive ones.
  \item Solvers are the same. Each resolves a part of the physics and models
        or ignores the rest.
  \item \pkgname{} keeps them behind one workflow, so the choice of fidelity
        is a choice of workspace and not a change of habits.
\end{itemize}
\begin{center}
\begin{tikzpicture}[
  s/.style={draw=hanblue,fill=hanblue!8,rounded corners=2pt,minimum height=11mm,
            text width=18mm,align=center,font=\scriptsize,inner sep=1.5pt},
  d/.style={s,dashed}]
  \node[s] (x) at (0,0) {\textbf{XFOIL}\\2D panel and boundary layer};
  \node[s] (q) at (2.3,0) {\textbf{QPROP}\\blade element};
  \node[d] (f) at (4.6,0) {\textbf{FlightStream}\\3D panel method};
  \node[s] (s) at (6.9,0) {\textbf{SU2}\\Euler, RANS};
  \node[s] (o) at (9.2,0) {\textbf{OpenFOAM}\\RANS};
  \node[s] (p) at (11.5,0) {\textbf{PyFR}\\high order, time resolved};
  \draw[-{Stealth[length=6pt]},thick,SoftGray] (-1.0,-0.95) -- (12.5,-0.95)
    node[midway,below,font=\scriptsize,text=SoftGray]
    {more of the flow resolved, more cost per point};
\end{tikzpicture}
\end{center}
{\scriptsize\color{SoftGray}Dashed: run through pyflightstream, in its own
workspace format.\par}
\end{frame}

\begin{frame}{What each solver answers}
\relax{\cfdbtablesize
\begin{tabularx}{\linewidth}{@{}P{0.13\linewidth}P{0.27\linewidth}P{0.27\linewidth}R@{}}
\toprule
\textbf{Solver} & \textbf{It models} & \textbf{One point is} & \textbf{It
returns} \\
\midrule
XFOIL [5] & a 2D airfoil: panels coupled to an integral boundary layer & one
  angle of attack & \code{CL}, \code{CD}, \code{CDp}, \code{CM} \\
QPROP [6] & an isolated propeller in axial flow, from blade sections and
  their 2D polars & one speed and rotation rate & thrust, torque, power,
  \code{CT}, \code{CP}, efficiency \\
SU2 [7] & compressible flow on a volume mesh, Euler or RANS & one Mach number
  and attitude & force and moment coefficients \\
OpenFOAM [8] & steady incompressible RANS, \code{simpleFoam}, on a 2D
  extruded mesh & one speed and angle of attack & \code{Cd}, \code{Cl},
  \code{CmPitch} \\
PyFR [9] & 2D compressible Navier-Stokes, high order, marched in time & one
  Mach, Reynolds number and angle & time-averaged \code{Cl}, \code{Cd} \\
\bottomrule
\end{tabularx}}
\vspace{1mm}
{\small\color{SoftGray}Every library solver runs \textbf{steady} points
today. Guide 3 states the limits of each solver.\par}
\end{frame}

\begin{frame}{The mock solver}
\begin{columns}[T]
\begin{column}{0.5\textwidth}
\begin{itemize}
  \item \code{mock} is a solver that computes nothing and returns fixed
        numbers.
  \item It needs no installation, no mesh and no licence.
  \item Its setup file can ask for a converged point, a point at the
        iteration limit or a failed one.
\end{itemize}
\end{column}
\begin{column}{0.46\textwidth}
\begin{block}{Use it to}
\begin{itemize}\small
  \item learn the commands (Part 4 of this guide);
  \item try a matrix layout before the meshes exist;
  \item see what a failed point looks like in the records, safely.
\end{itemize}
\end{block}
\end{column}
\end{columns}
\begin{takeaway}Start on the mock. Move to your solver when the workflow is
no longer the thing you are learning.\end{takeaway}
\end{frame}

\begin{frame}{A result is evidence, not validation}
\begin{itemize}
  \item \pkgname{} records what was run, with which inputs, and what the
        solver reported. It keeps the hashes of the inputs beside each
        result.
  \item It classifies each point: \code{converged}, \code{iteration_limit}
        or \code{failed}. What ``converged'' rests on is the solver's own
        evidence, stated in guide 3.
  \item It does \textbf{not} judge the mesh, the turbulence model or the
        physics. A converged point on a poor mesh is still a converged
        point.
\end{itemize}
\begin{cfdbevidence}[What you can rely on]
The same matrix row gives the same input deck, byte for byte, and the record
names every file the point read.
\end{cfdbevidence}
\end{frame}

% Copyright (c) 2026 Geovana Neves. Licensed under CC BY 4.0 (https://creativecommons.org/licenses/by/4.0/). See LICENSE-AND-AUTHORSHIP.md.
%
% overview/text/02-workflow-and-first-campaign.tex
%
% Parts 3 to 5 of guide 1: the workflow in one picture, a first campaign on
% the mock solver, and the map of the guides.

\section{The workflow}

\begin{frame}{Seven stages, one direction}
\begin{center}
\begin{tikzpicture}[
  st/.style={draw=hanblue,fill=hanblue!8,rounded corners=2pt,minimum height=9mm,
             text width=17mm,align=center,font=\scriptsize,inner sep=2pt},
  ar/.style={-{Stealth[length=5pt]},thick,hanblue}]
  \node[st] (a) at (0,0) {\textbf{1. Create}\\\code{init}, inputs};
  \node[st] (b) at (2.5,0) {\textbf{2. Plan}\\\code{plan}};
  \node[st] (c) at (5,0) {\textbf{3. Run}\\\code{run}};
  \node[st] (d) at (7.5,0) {\textbf{4. Follow}\\\code{status}, \code{log}};
  \node[st] (e) at (10,0) {\textbf{5. Collect}\\\code{collect}, \code{post}};
  \node[st,text width=30mm] (f) at (7.5,-1.6) {\textbf{6. Correct the records}\\\code{mark-failed}, \code{restore}};
  \node[st,text width=30mm] (g) at (10.9,-1.6) {\textbf{7. Manage storage}\\\code{free-space}, \code{sync}};
  \draw[ar] (a) -- (b); \draw[ar] (b) -- (c); \draw[ar] (c) -- (d); \draw[ar] (d) -- (e);
  \draw[ar,SoftGray,dashed] (d.south) -- (f.north);
  \draw[ar,SoftGray,dashed] (e.south) -- (e.south|-f.north);
\end{tikzpicture}
\end{center}
\begin{itemize}\small
  \item Stages 1 to 5 are every campaign. Stages 6 and 7 are for the day a
        record is wrong or the disk is full.
  \item \code{cfd-workspace} does stage 1. \code{cfd-matrix} does the rest.
  \item The cheatsheet, guide 4, has one page per stage.
\end{itemize}
\end{frame}

\begin{frame}[fragile]{The run matrix: one row per simulation}
\begin{lstlisting}
POL  | RUN | FLIGHT_CONDITION                  | SWEEP_VALUES | GEOMETRY | REF | SET | PPROC
-----|-----|-----------------------------------|--------------|----------|-----|-----|------
5001 | 1   | MACH:0.85, REmi:5, ALPHA:sweep    | 2.5,2.75     | wb.su2   | 001 | 001 | 001
5002 | 0   | MACH:0.85, REmi:5, ALPHA:3.0      | -            | wb.su2   | 001 | 002 | 001
\end{lstlisting}
\begin{itemize}\small
  \item \code{inputs/matrices/<name>.fs}: text, columns separated by
        \code{|}. A cell holding \code{-} states nothing.
  \item \code{POL} names the simulation. \code{RUN} 1 runs the row, 0 keeps
        it in the file without running it.
  \item \code{FLIGHT_CONDITION} holds \code{KEY:value} pairs. One key may say
        \code{sweep}; its values are in \code{SWEEP_VALUES}, one
        \textbf{point} each.
  \item \code{SET}, \code{REF} and \code{PPROC} name the three input files
        of the row.
\end{itemize}
\end{frame}

\begin{frame}{Three input files per row: s, r, p}
\relax{\cfdbbodysize
\begin{tabularx}{\linewidth}{@{}P{0.12\linewidth}P{0.38\linewidth}R@{}}
\toprule
\textbf{Column} & \textbf{File} & \textbf{It holds} \\
\midrule
\code{SET} & \code{inputs/setups/s<id>.toml} & how the solver runs: model,
  numerics, iteration limit. The keys belong to the solver. \\
\code{REF} & \code{inputs/references/r<id>.toml} & reference area, reference
  length and moment origin. \\
\code{PPROC} & \code{inputs/pproc/p<id>.toml} & what to keep from each run
  beyond the coefficients. \\
\code{GEOMETRY} & \code{inputs/mesh/<file>} & the mesh or geometry, in the
  format the solver reads. \\
\bottomrule
\end{tabularx}}
\vspace{2mm}
\begin{itemize}\small
  \item \code{001} and \code{s001} in the \code{SET} cell both name
        \code{s001.toml}.
  \item Many rows may share one file. A value the row and a file both state
        must agree.
\end{itemize}
\begin{takeaway}You write only under \code{inputs/}. Nothing a command writes
goes there.\end{takeaway}
\end{frame}

\begin{frame}{Plan first, then run}
\begin{columns}[T]
\begin{column}{0.48\textwidth}
\begin{block}{\code{cfd-matrix plan <matrix>}}
\begin{itemize}\small
  \item Reads the matrix and every file it names.
  \item Lists each point and the files its run would write.
  \item Launches nothing.
  \item Saves the plan receipt, \code{post/<matrix>/plan.json}.
\end{itemize}
\end{block}
\end{column}
\begin{column}{0.48\textwidth}
\begin{block}{\code{cfd-matrix run <matrix>}}
\begin{itemize}\small
  \item Needs the receipt.
  \item Refuses when an input changed after the plan: plan again.
  \item Writes one folder per point under \code{sims/} and one record per
        point in \code{runs.json}.
\end{itemize}
\end{block}
\end{column}
\end{columns}
\begin{takeaway}A mistake in an input is found by \code{plan}, in seconds,
and not by a job that waited in a queue.\end{takeaway}
\end{frame}

\begin{frame}{Records, and nothing deleted}
\begin{itemize}
  \item \code{runs.json} holds one record per point: its inputs and their
        hashes, its status, its times and the verdict on its convergence.
  \item A point that is already recorded \textbf{refuses} a second run.
        \code{--resume} runs only what is not recorded yet, and says how to
        redo a failed point; \code{--force-rerun} redoes a point you name.
  \item Before a record or a result is replaced, the old one moves to
        \code{archive/<stamp>/}.
  \item The commands that correct records or manage storage only
        \textbf{preview} until you add \code{--apply}.
\end{itemize}
\begin{takeaway}Finished evidence is never replaced by accident.\end{takeaway}
\end{frame}

\begin{frame}{Three verdicts}
\relax{\cfdbbodysize
\begin{tabularx}{\linewidth}{@{}P{0.22\linewidth}R@{}}
\toprule
\textbf{Verdict} & \textbf{Meaning} \\
\midrule
\code{converged} & the solver ended and its own evidence of convergence is
  there \\
\code{iteration_limit} & the solver ended properly without that evidence;
  the numbers exist and are not settled \\
\code{failed} & the process failed, or its output is missing, inconsistent
  or not a number \\
\bottomrule
\end{tabularx}}
\vspace{2mm}
\begin{itemize}\small
  \item The verdict is beside every row of the collected table.
  \item You may overrule it, with a reason on record:
        \code{cfd-matrix mark-failed} and \code{cfd-matrix mark-converged}.
\end{itemize}
\end{frame}

\section{A first campaign}

\begin{frame}[fragile]{Four small files}
\begin{lstlisting}
# inputs/matrices/first.fs
POL  | RUN | FLIGHT_CONDITION      | SWEEP_VALUES | REF | SET | PPROC
1001 | 1   | MACH:0.2, ALPHA:sweep | 0,2,4        | 001 | 001 | 001
\end{lstlisting}
\begin{columns}[T]
\begin{column}{0.36\textwidth}
\begin{lstlisting}[style=cfdbtoml]
# inputs/setups/s001.toml
[mock]
outcome = "converged"
iterations = 10
\end{lstlisting}
\begin{lstlisting}[style=cfdbtoml]
# inputs/pproc/p001.toml
history = true
\end{lstlisting}
\end{column}
\begin{column}{0.6\textwidth}
\begin{lstlisting}[style=cfdbtoml]
# inputs/references/r001.toml
area = 1.5
length = 1.0
moment_origin = [0, 0, 0]
axis_system = "body"
axis_directions = ["aft", "right", "up"]
\end{lstlisting}
\end{column}
\end{columns}
\vspace{1mm}
{\small One simulation, \code{1001}, with three points: angle of attack 0, 2
and 4 degrees at Mach 0.2.\par}
\end{frame}

\begin{frame}[fragile]{Six commands}
\begin{lstlisting}
cfd-workspace init campaign --solver mock
cd campaign
   (write the four files of the previous slide)
cfd-matrix plan first
cfd-matrix run first
cfd-matrix status --points
cfd-matrix collect first
\end{lstlisting}
\begin{itemize}\small
  \item \code{plan} lists three points. \code{run} runs them on this machine.
  \item \code{status --points} shows each point with its status and verdict.
  \item \code{collect} writes \code{post/first/campaign_sweep.csv}: one row
        per point, with the flight condition, the air state, the verdict and
        the coefficients.
\end{itemize}
\begin{takeaway}Change \code{outcome} to \code{"failed"} in a second setup
file and run a second row: you have seen a failure before it matters.\end{takeaway}
\end{frame}

\begin{frame}{From the mock to your solver}
\begin{enumerate}
  \item Make a new workspace: \code{cfd-workspace init <folder> --solver
        <name>}.
  \item Put the mesh in \code{inputs/mesh/} and name it in the
        \code{GEOMETRY} cell.
  \item Write the setup file with the solver's keys: guide 3 has one worked
        setup per solver.
  \item Say where the solver is on this machine: in the setup's
        \code{[solver]} table, or in \code{inputs/executables.toml} with the
        \code{SOLVER_BUILD} column. A run refuses before it records anything
        when the solver is missing.
  \item Plan, run, follow, collect: the same commands as on the mock.
\end{enumerate}
\end{frame}

\section{The guides}

\begin{frame}{The four guides}
\relax{\cfdbbodysize
\begin{tabularx}{\linewidth}{@{}P{0.05\linewidth}P{0.37\linewidth}R@{}}
\toprule
 & \textbf{File} & \textbf{Read it when} \\
\midrule
1 & \code{cfdb-guide-01-overview.pdf} & first: the idea, the solvers, the
  workflow \\
2 & \code{cfdb-guide-02-workspaces.pdf} & you start a campaign: the tree,
  the inputs, every stage, a workstation and a cluster \\
3 & \code{cfdb-guide-03-solvers.pdf} & you write a setup file or read a
  result: one part per solver \\
4 & \code{cfdb-guide-04-cheatsheet.pdf} & you are at the keyboard: every
  command and option, one page per stage \\
\bottomrule
\end{tabularx}}
\vspace{2mm}
\begin{cfdbevidence}[FlightStream]
A FlightStream workspace is run by \pyfsname{} [4]. Its commands, inputs and
guides are those of \pyfsname: \url{\pyfssite}.
\end{cfdbevidence}
\end{frame}

% Geometry, mesh and result products in the current workflow.
\section{Geometry, mesh and results}

\begin{frame}[fragile]{Build the shape, then the mesh}
\begin{lstlisting}
cfd-geom validate model.toml
cfd-geom run model.toml
cfd-mesh airfoil-2d --naca 0012 --reynolds 6000000
\end{lstlisting}
\begin{itemize}\small
  \item \code{cfd-geom} builds parametric geometry and records its stages in
        \code{inputs/geometry/<name>/}, with \code{geometry.json}.
  \item \code{cfd-mesh} generates and transforms meshes in \code{inputs/mesh/},
        with provenance sidecars. The airfoil command writes SU2 and gmsh files.
  \item A 2D SU2 or OpenFOAM row that sweeps \code{ALPHA} can become a section
        polar: \code{polars/P<POL>.txt}, readable by XFOIL and QPROP.
\end{itemize}
\end{frame}

\begin{frame}{From a saved solution to a recorded result}
\begin{itemize}\small
  \item \code{RESTART} in a row names a recorded source solution;
        \code{run --restart} continues each point from its own solution.
        SU2, OpenFOAM and PyFR support both. Earlier evidence is archived first.
  \item The host guard is \textbf{off by default}. Add
        \code{--one-run-at-a-time}, or set \code{one_run_at_a_time = true}
        under \code{[solver]} in \code{init.toml}, to reserve the host.
  \item In \code{VAR_NAMES_VALUES}, \code{hpc: alpha / cluster: west}
        selects a profile and cluster per row. \code{hpc:} overrides
        \code{--hpc}; \code{--local} ignores both row selections.
  \item \code{collect} writes \code{campaign_summary.json}, family tables such as
        \code{FAM-ACOEF-P<POL>.csv}, group tables and \code{products.json}.
        Setup \code{[groups]} boundary forces stay in the summary, separately
        from the family-based group CSV files. Guide 2 gives the file map.
\end{itemize}
\end{frame}


\begin{frame}{References}
\begin{reflist}
  \bibitem{g1r01} \docref{workspace-layout}{Workspace layout}
  \bibitem{g1r02} \docref{cli}{Command line}
  \bibitem{g1r03} \docref{matrix-columns}{Matrix columns}
  \bibitem{g1r04} \docref{flightstream-workspaces}{FlightStream workspaces}
  \bibitem{g1r05} M. Drela, ``XFOIL: an analysis and design system for low
        Reynolds number airfoils'', in \emph{Low Reynolds Number
        Aerodynamics}, Lecture Notes in Engineering 54, Springer, 1989.
  \bibitem{g1r06} M. Drela, ``QPROP formulation'', MIT Aero and Astro, 2006.
  \bibitem{g1r07} T. D. Economon, F. Palacios, S. R. Copeland, T. W. Lukaczyk
        and J. J. Alonso, ``SU2: an open-source suite for multiphysics
        simulation and design'', \emph{AIAA Journal}, 54(3), pp. 828-846,
        2016.
  \bibitem{g1r08} H. G. Weller, G. Tabor, H. Jasak and C. Fureby, ``A
        tensorial approach to computational continuum mechanics using
        object-oriented techniques'', \emph{Computers in Physics}, 12(6),
        pp. 620-631, 1998.
  \bibitem{g1r09} F. D. Witherden, A. M. Farrington and P. E. Vincent,
        ``PyFR: an open source framework for solving advection-diffusion type
        problems on streaming architectures using the flux reconstruction
        approach'', \emph{Computer Physics Communications}, 185(11),
        pp. 3028-3040, 2014.
\end{reflist}
\end{frame}

\end{document}
